\section{Related work}\label{related-work}

\subsection{Encoder--decoder architectures for medical image segmentation}

Fully convolutional encoder--decoder networks remain a central design pattern
for biomedical image segmentation. U-Net introduced a symmetric contracting
and expanding architecture in which high-resolution encoder features are
transferred to the decoder through skip connections, allowing semantic context
and spatial detail to be combined during prediction \cite{unet}. UNet++ further
showed that the structure of skip pathways is important by using nested dense
connections to reduce the semantic discrepancy between encoder and decoder
representations \cite{unetpp}. These works established the importance of
hierarchical feature reuse, but they do not prescribe how the decoder should
adapt its receptive-field mixture to the image-specific content of a polyp.

Polyp-specific convolutional models have explored efficient alternatives to
large attention-based designs. HarDNet-MSEG combines a low-memory-traffic
convolutional backbone with a lightweight decoder and reports strong
polyp segmentation can be obtained without a transformer-heavy architecture
\cite{hardnetmseg}. PraNet introduced a parallel partial decoder and reverse
attention to progressively refine uncertain polyp regions \cite{pranet}. These
methods highlight the value of efficient decoding and iterative refinement, but
their main mechanisms are not designed to use frequency information from the
bottleneck to regulate multi-scale convolutional branches at an earlier decoder
stage.

Modern convolutional encoders provide another route to high-capacity dense
prediction. ConvNeXt revisits convolutional network design with large-kernel
depthwise convolution, inverted channel expansion, Layer Normalization, GELU
activation, layer scaling, and stochastic depth, producing a hierarchical
backbone with strong representation capacity \cite{convnext}. We use
ConvNeXt-Tiny as the encoder and focus the new mechanism on the interaction
between bottleneck frequency information and Stage-3 multi-scale convolution.

\subsection{Transformer and contextual polyp segmentation}

The large variation in polyp size, shape, color, and boundary contrast motivates
architectures that capture both local appearance and broader anatomical context.
Transformer-based approaches address this need by modeling long-range spatial
relationships. TransUNet combines convolutional feature extraction with a
transformer encoder for medical image segmentation \cite{transunet}, and
Polyp-PVT applies pyramid vision transformer features and specialized fusion to
polyp segmentation \cite{polyppvt}. Recent polyp-segmentation studies have
continued this direction through contrastive transformer modeling, contextual
information flow, phase-wise feature pyramids, and pyramid-transformer feature
fusion \cite{ctnet,cifformer,pfprnet,mfnet}.

These approaches show that global context and cross-level interaction are useful
for polyp segmentation. However, attention-based or multi-stage contextual
architectures can introduce additional architectural complexity. The network
therefore keeps a convolutional encoder--decoder form and examines whether a
compact bottleneck frequency descriptor can guide the
relative weighting of efficient multi-scale convolutional branches in the first
decoder fusion stage.

\subsection{Efficient multi-scale convolutional decoding}

Multi-scale convolution is widely used because polyps may appear as small,
subtle lesions or occupy a large image region. Parallel kernels and dilated or
depthwise branches allow a decoder to aggregate local and broader responses
without relying entirely on self-attention. EMCAD introduced an efficient
multi-scale convolutional attention decoder for medical image segmentation and
uses Multi-Scale Convolution Blocks (MSCBs) as part of its lightweight decoding
design \cite{emcad}. In an MSCB, multiple convolutional branches provide
different receptive fields and are combined to form a multi-scale response.

MSCB serves as the multi-scale reference path. A conventional MSCB aggregation treats the branch
mixture as fixed by the block design rather than conditioning the branch weights
on the image-specific frequency characteristics extracted at the bottleneck. Our
RFG-MSCB retains the ordinary MSCB path as a reference and adds only a bounded,
frequency-conditioned residual deviation from that path.

\subsection{Frequency-aware and adaptive feature recalibration}

Attention and recalibration mechanisms provide a general precedent for using
compact descriptors to control feature responses. Squeeze-and-Excitation
Networks use global channel descriptors to recalibrate convolutional features
\cite{senet}, while Selective Kernel Networks adaptively weight convolutional
branches with different receptive fields \cite{sknet}. These methods are
important novelty comparators because they show that descriptor-driven
reweighting and adaptive receptive-field selection are established ideas.

Frequency-aware attention provides a closer conceptual link. FcaNet interprets
channel attention through frequency analysis and uses frequency components for
channel attention \cite{fcanet}. In polyp segmentation, Polyp-Mamba combines
multi-frequency perception with a gated selection mechanism to model
frequency-sensitive representations \cite{polypmamba}. These studies motivate
the usefulness of frequency information, but they do not correspond to the
specific coupling used here: a frequency descriptor produced by a bottleneck
FAFEM is transferred across levels and used to regulate Stage-3 MSCB branch
weights.

We adopt the Frequency-Aware Feature Enhancement Module (FAFEM) at the bottleneck
to enhance deep features and produce a low-/high-frequency descriptor. Its
descriptor is reused as cross-level guidance for multi-scale convolution in the
decoder.

\subsection{Boundary and recent polyp-segmentation refinements}

Weak and visually ambiguous lesion boundaries remain a major challenge in polyp
segmentation. Boundary-aware methods often introduce explicit edge supervision,
auxiliary contour branches, image-pyramid high-frequency cues, or reverse
attention. MEGANet employs multi-scale edge-guided attention for weak-boundary
polyp segmentation \cite{meganet}, and MSBP-Net performs multi-scale boundary
prediction to improve contour delineation \cite{msbpnet}. Other recent methods
model uncertainty, contextual flow, frequency-sensitive representations, and
multi-stage refinement \cite{mein,cifformer,polypmamba,pranetv2}.

The network does not add an explicit boundary branch or edge label; it uses the
boundary-aware loss described in the Experiments section. Boundary-specific
improvement is not evaluated separately, so the frequency-guided decoder should
not be described as a dedicated boundary module.

\subsection{Positioning of the method}

Existing work has improved polyp segmentation through U-shaped skip reuse,
reverse attention, transformer-based context modeling, efficient convolutional
decoding, frequency-aware representation learning, and boundary-aware
refinement. The method is positioned at the intersection of efficient
multi-scale convolution and frequency-aware recalibration.

Specifically, the final model uses a ConvNeXt-Tiny encoder, a bottleneck FAFEM,
a lightweight residual bottleneck, Stage-3 SimpleFusion, the Residual
Frequency-Guided MSCB, and standard remaining decoder stages. We adopt FAFEM and
the MSCB structure; the new component is their coupling: the bottleneck FAFEM
frequency descriptor adaptively weights MSCB branches
at Stage 3, the guided response is introduced as a bounded residual deviation
from ordinary MSCB aggregation, and the guidance forms a bottleneck-to-Stage-3
cross-level connection. This separates the RFG-MSCB mechanism from the components
on which it is built.
